% !TeX root = ../main.tex
\begin{frame}{课程评价——李牧琛}
  \begin{SlideCanvas}
    \node[anchor=north west,inner sep=0pt,text width=850pt,align=left,
      text=BICNavy] at (55,80) {\begin{minipage}{850pt}
      \usebeamerfont{normal text}\raggedright
      \setlength{\parindent}{0pt}\setlength{\parskip}{5pt}
\par\addvspace{8pt}{\usebeamerfont{BIC point}一、课程评价\par}\nobreak
本课程以专题讲座形式展开，覆盖医疗传感与器件、医学探测与成像、临床诊疗和智能机器人等方向。课程使我更加关注临床需求、技术原理和应用条件之间的联系。\par
\par\addvspace{8pt}{\usebeamerfont{BIC point}二、学习收获\par}\nobreak
声学软材料、柔性电子、MEMS超声及影像探测等专题，帮助我从材料、结构和信号层面理解医学任务。结合本次血管支架汇报，我认识到应将病变覆盖、管腔支撑等需求转成可检验指标，并区分单胞仿真、器件试验和临床效果的证据层次。\par
\par\addvspace{8pt}{\usebeamerfont{BIC point}三、意见与建议\par}\nobreak
1. 课前提供术语与基础导读，帮助不同专业学生理解临床问题。\par
2. 增加器件、样机或仿真演示，说明系统构成与使用条件。\par
3. 围绕同一临床问题组织跨专业讨论，贯通需求、设计和验证。\par
    \end{minipage}};
  \end{SlideCanvas}
  \note{本页为按用户提供的十二个专题讲座及本次支架汇报背景撰写的李牧琛课程评价。李牧琛《医工交叉研究前沿进展》课程反馈与建议。一、课程评价\par 本课程以专题讲座形式展开，覆盖医学探测与成像、医疗传感与柔性器件、临床诊疗和智能机器人等方向。不同专题将技术原理与医学应用联系起来，使我能够从多学科视角理解研究问题。课程带给我的启发，是从临床需求出发选择工程方案，并同时思考技术的适用条件与验证方式。\par \par 二、学习收获\par 声学软材料、柔性电子和MEMS超声换能器等专题，让我更加关注材料、器件结构与使用条件的关系；弱光探测、磁共振影像和人工听觉等内容，则帮助我理解感知器件、信号与医学任务之间的联系。\par \par 神经外科、消化道肿瘤内镜切除和胃肠道转化医学等讲座，使我注意到临床操作与应用场景对工程设计的约束。软体智能、人机共演、超声医学中的机器人及人工智能、智能手术机器人等专题，进一步提示我重视系统集成和操作者的实际使用。\par \par 结合本次血管支架汇报，我更加明确：医工交叉研究需要把病变覆盖、管腔支撑等需求转化为长度变化、径向力和局部应变等可检验指标，并区分单胞仿真、器件试验与临床效果的证据层次。这也将影响我今后设置研究目标、边界条件和验证方案的方式。\par \par 三、意见与建议\par 1. 课前提供术语、基础原理和临床需求导读，帮助不同专业背景的学生建立共同知识基础。\par \par 2. 增加器件、样机或仿真演示，结合超声、内镜、柔性传感和机器人等实例说明系统构成与使用条件。\par \par 3. 围绕具体临床问题组织跨专业讨论，串联需求、设计指标和验证路径，帮助学生把各专题的启发落实为研究思路。}
\end{frame}
